(1) 绳速 \(v=\dot x=R\dot\varphi\)，滑轮转动动能为 \(Iv^2/(2R^2)\)。将两个质点与滑轮整体考虑，绳张力为内力，故有效惯性 \(A=m_1+m_2+I/R^2\)、有效线性阻尼 \(B=b/R^2\)，方程严格为 \(A\ddot x+B\dot x=(m_1-m_2)g\equiv F\)，初态速度与位移另行给定。
(2) 对 \(b>0\)，令 \(\tau=A/B\) 和 \(v_\infty=F/B\)。线性初值问题 \(A\dot v+Bv=F\) 且 \(v(0)=0\) 给出 \(v(t)=v_\infty(1-e^{-t/\tau})\)；再次积分得 \(x(t)-x(0)=v_\infty[t-\tau(1-e^{-t/\tau})]\)，它在 \(t=0\) 的速度与加速度都满足原方程。
(3) 乘以速度，\(Fv=Av\dot v+Bv^2\)。第一项是 \(d(Av^2/2)/dt\)，第二项 \(Bv^2=b\dot\varphi^2\ge0\) 是轴承的耗散功率。\([A]=\mathrm{kg}\)、\([B]=\mathrm{kg\,s^{-1}}\)，故每一项均为功率，且阻力不能凭空增加机械能。
(4) 终端速度为 \(F/B=FR^2/b\)；在 \(t\) 固定时作指数展开 \(1-e^{-Bt/A}=Bt/A+O(B^2)\)，得到 \(v(t)\to Ft/A\)、\(x(t)-x(0)\to Ft^2/(2A)\)。终端极限 \(t\to\infty\) 与无阻尼极限 \(b\to0\) 不可交换，因为后者没有有限终端速度。
